\chapter{Dedica delle offerte}

%\suttaref{Trad.}%
[Yo so] bha꜕gavā a꜕rahaṁ sammāsambuddho

\begin{english}
[Al Beato] al Nobile e Perfettamente Risvegliato
\end{english}

Svākkhā꜓to yena bha꜕gava꜓tā dhammo

\begin{english}
alla Verità, che ha così ben esposto
\end{english}

Supaṭi꜕panno yassa bha꜕gava꜕to sāvaka꜕saṅgho

\begin{english}
ai discepoli del Beato, che hanno ben praticato
\end{english}

Tam-ma꜓yaṁ bha꜕gavantaṁ sa꜕dhammaṁ sa꜕saṅghaṁ

\begin{english}
al Buddha, al Dhamma e al Sangha
\end{english}

Imehi꜓ sakkārehi꜕ yathārahaṁ āropi꜕tehi a꜕bhi꜓pūja꜕yāma

\begin{english}
offriamo i nostri onori e i nostri omaggi.
\end{english}

Sādhu꜓ no bhante bha꜕gavā su꜕cira-parinibbu꜕topi

\begin{english}
è bene per noi che il Beato, pur avendo ottenuto la Liberazione
\end{english}

Pacchi꜓mā-ja꜕na꜓tānu꜓kampa꜕-mānasā

\begin{english}
abbia avuto compassione, delle generazioni future.
\end{english}

Ime sakkāre dugga꜕ta꜕-paṇṇākāra꜓-bhūte pa꜕ṭiggaṇhātu

\begin{english}
questi umili onori e doni siano accolti
\end{english}

Amhā꜓kaṁ dīgha꜕rattaṁ hi꜕tāya su꜕khāya

\begin{english}
a nostro duraturo beneficio, e per la felicità che questo ci dona.
\end{english}

\clearpage

Arahaṁ sammāsambuddho bha꜕gavā

\begin{english}
il Beato, il Nobile Perfettamente Risvegliato
\end{english}

Buddhaṁ bha꜕gavantaṁ a꜕bhi꜓vādemi

\begin{english}
rendo omaggio al Buddha, al Beato.
  \instr{Inchino}
\end{english}

[Svākkhā꜓to] bha꜕gava꜓tā dhammo

\begin{english}
[L'Insegnamento] completamente esposto dal Beato
\end{english}

Dhammaṁ namassāmi

\begin{english}
mi inchino al Dhamma.
  \instr{Inchino}
\end{english}

[Supaṭi꜕panno] bha꜕gava꜕to sāvaka꜕saṅgho

\begin{english}
[I discepoli del Beato] che hanno ben praticato
\end{english}

Sa꜓ṅghaṁ na꜕māmi

\begin{english}
mi inchino al Sangha.
  \instr{Inchino}
\end{english}

\chapter{Omaggio Preliminare}

\begin{leader}
  [Ha꜓nda mayaṁ buddhassa꜕ bha꜕gavato pubbabhāga-namakā꜕raṁ karomase]
\end{leader}

\begin{english}
[Rendiamo ora omaggio preliminare al Buddha]
\end{english}

\vspace{\baselineskip}

%\suttaref{Cf. DN 21}%
Namo tassa bha꜕gava꜕to araha꜕to sa꜓mmāsa꜓mbuddha꜕ssa

\instr{Three times}

\begin{english}
Omaggio al Beato, al Nobile Perfettamente Risvegliato. (3 volte)

  \instr{Three times}
\end{english}

\clearpage

\chapter{Omaggio al Buddha}

\begin{leader}
  [Ha꜓nda mayaṁ buddhābhi꜕tthu꜕tiṁ karomase]
\end{leader}

\begin{english}
[Cantiamo ora la lode del Buddha]
\end{english}

%\suttaref{DN 2}%
Yo so tathā꜓ga꜕to a꜕rahaṁ sammāsambuddho

\begin{english}
il Tathagata è il Puro, Perfettamente Risvegliato
\end{english}

Vijjāca꜕raṇa꜓-sampanno

\begin{english}
impeccabile in condotta e conoscenza
\end{english}

Su꜕ga꜕to

\begin{english}
il Sereno,
\end{english}

Loka꜕vi꜓dū

\begin{english}
conoscitore dei mondi
\end{english}

Anu꜓tta꜕ro purisa꜕damma-sārathi

\begin{english}
abile insegnante per chi cerca educazione
\end{english}

Satthā deva-ma꜕nussānaṁ

\begin{english}
maestro di déi e di uomini
\end{english}

Buddho bha꜕gavā

\begin{english}
Risvegliato e Beato,
\end{english}

Yo imaṁ lokaṁ sa꜕devakaṁ sa꜕mārakaṁ sa꜕brahma꜕kaṁ

\begin{english}
ha rivelato la Verità, con profonda comprensione CHECK ORDER HERE
\end{english}

Sassa꜓maṇa-brāhmaṇiṁ pa꜕jaṁ sa꜕deva-ma꜕nussa꜓ṁ sa꜕yaṁ a꜕bhiññā sacchika꜕tv꜓ā pa꜕vedesi

\begin{english}
in questo mondo con i suoi déi, demoni e Brahma, i suoi ricercatori e saggi, le guide e loro popoli.
\end{english}

Yo dhammaṁ dese꜓si ā꜕di꜓-kalyāṇaṁ majjhe꜓-ka꜕lyāṇaṁ \\pa꜕riyosāna-k꜕alyāṇaṁ

\begin{english}
lui ha indicato il Dhamma, bello all'inizio, nel mezzo e alla fine.
\end{english}

Sāttha꜓ṁ sa꜕byañjanaṁ kevala-pa꜕ripuṇṇaṁ pa꜕risuddhaṁ \\brahma-ca꜕ri꜓yaṁ pa꜕kāsesi

\begin{english}
ha insegnato la Vita Santa, completa e pura in lettera ed essenza
\end{english}

Tam-aha꜓ṁ bha꜕gavantaṁ a꜕bhi꜓pūja꜕yāmi tam-aha꜓ṁ bha꜕gavantaṁ \\si꜕rasā꜓ na꜕māmi

\begin{english}
rendo lode al Beato, mi inchino al Beato. 
  \instr{Inchino}
\end{english}

\clearpage

\chapter{Omaggio al Dhamma}

\begin{leader}
  [Ha꜓nda mayaṁ dhammābhi꜕tthu꜕tiṁ karomase]
\end{leader}

\begin{english}
[Cantiamo ora la lode del Dhamma]
\end{english}

%\suttaref{DN 16}%
Yo so svākkhā꜓to bha꜕gava꜓tā dhammo

\begin{english}
il Dhamma è ben esposto dal Beato
\end{english}

Sa꜓ndiṭṭhi꜕ko

\begin{english}
presente qui ed ora, 
\end{english}

A꜕kāli꜕ko

\begin{english}
senza-tempo, 
\end{english}

Ehi꜕passi꜕ko

\begin{english}
incoraggia la verifica diretta
\end{english}

Opanayi꜕ko

\begin{english}
conduce oltre, 
\end{english}

Pa꜕cca꜕ttaṁ vedi꜓ta꜕bbo viññūhi

\begin{english}
va vissuto di persona da ciascun ricercatore.
\end{english}

Tam-aha꜓ṁ dhammaṁ a꜕bhi꜓pūja꜕yāmi tam-aha꜓ṁ dhammaṁ \\si꜕rasā꜓ na꜕māmi

\begin{english}
rendo lode all'Insegnamento, mi inchino alla Verità. 
  \instr{Inchino}
\end{english}

\clearpage

\chapter{Omaggio al Sangha}

\begin{leader}
  [Ha꜓nda mayaṁ saṅghābhi꜕tthu꜕tiṁ karomase]
\end{leader}

\begin{english}
[Cantiamo ora la lode del Sangha]
\end{english}

%\suttaref{DN 16}%
Yo so supaṭi꜕panno bha꜕gava꜕to sāvaka꜕saṅgho

\begin{english}
i discepoli del Beato che hanno ben praticato
\end{english}

Ujupaṭi꜕panno bha꜕gava꜕to sāvaka꜕saṅgho

\begin{english}
vivono rettamente
\end{english}

Ñāyapaṭi꜕panno bha꜕gava꜕to sāvaka꜕saṅgho

\begin{english}
vivono sapientemente
\end{english}

Sā꜓mīci꜕pa꜕ṭi꜕panno bha꜕gava꜕to sāvaka꜕saṅgho

\begin{english}
vivono con integrità
\end{english}

Yadidaṁ cattāri purisa꜕yugāni aṭṭha꜓ purisa꜕pugga꜕lā

\begin{english}
sono le quattro coppie, le otto specie di Nobili persone
\end{english}

Esa bha꜕gava꜕to sāvaka꜕saṅgho

\begin{english}
tali sono i discepoli del Beato
\end{english}

Āhu꜕neyyo

\begin{english}
sono degni di doni,
\end{english}

Pāhu꜕neyyo

\begin{english}
degni di ospitalità,
\end{english}

\clearpage % this keeps the pali and english together on the same page

Dakkhi꜕ṇeyyo

\begin{english}
degni di offerte, 
\end{english}

Añja꜕li-ka꜕ra꜓ṇīyo

\begin{english}
degni di rispetto
\end{english}

Anu꜓tta꜕raṁ puññakkhe꜕ttaṁ lokassa

\begin{english}
il più fertile terreno di meriti per il mondo. 
\end{english}

Tam-aha꜓ṁ saṅghaṁ a꜕bhi꜓pūja꜕yāmi tam-aha꜓ṁ saṅghaṁ \\si꜕rasā꜓ na꜕māmi

\begin{english}
rendo lode alle Nobili Persone, mi inchino al Sangha. 
  \instr{Inchino}
\end{english}

\clearpage

\chapter{Omaggio alla triplice Gemma}

\begin{leader}
  [Ha꜓nda mayaṁ ratanattaya-paṇāma-gāthā꜓yo c'eva\\
  sa꜓ṁvega-parikittana-pāṭhañca꜕ bhaṇāmase]
\end{leader}

\begin{english}
  [Cantiamo ora la lode della Triplice Gemma e un passaggio che susciti urgenza]
\end{english}

%\suttaref{Trad.}%
Buddho su꜕suddho ka꜕ruṇā-maha꜓ṇṇavo

\begin{english}
il Buddha, assolutamente puro, oceano di compassione
\end{english}

Yo'ccanta꜕-suddhabba꜕ra-ñāṇa꜕-loca꜕no

\begin{english}
colui che possiede l'occhio della saggezza purissima e sublime
\end{english}

Lokassa꜕ pāpūpa꜕ki꜓lesa꜕-ghāta꜕ko

\begin{english}
distruttore della corruzione interiore
\end{english}

Vandāmi꜓ buddhaṁ a꜕ha꜓m-āda꜕rena꜕ taṁ

\begin{english}
con devozione rendo omaggio al Buddha.
\end{english}

Dhammo pa꜕dīpo vi꜕ya tassa꜕ satthu꜕no

\begin{english}
il Dhamma, come una lampada,
\end{english}

Yo magga꜓-pākāma꜕ta꜕-bheda꜕-bhinna꜕ko

\begin{english}
illumina il sentiero e il suo frutto: il Senza-Morte,
\end{english}

Lokuttaro yo ca꜕ ta꜕d-attha꜕-dīpa꜕no

\begin{english}
ciò che è al di là del mondo condizionato
\end{english}

Vandāmi꜓ dhammaṁ a꜕ha꜓m-āda꜕rena꜕ taṁ

\begin{english}
con devozione rendo omaggio al Dhamma. 
\end{english}

\clearpage

Sa꜓ṅgho su꜕khettābhyati-khe꜕tta-sa꜓ññito

\begin{english}
il Sangha, eccellente campo di meriti, disciplinato e contenuto
\end{english}

Yo diṭṭha꜓-santo su꜕ga꜕tānu꜕bodha꜕ko

\begin{english}
ha realizzato la Pace, composto dai Risvegliati dopo il Beato
\end{english}

Lolappa꜕hīno a꜕ri꜓yo su꜕medha꜕so

\begin{english}
nobili e saggi, liberi da ogni brama
\end{english}

Vandāmi꜓ saṅghaṁ a꜕ha꜓m-āda꜕rena꜕ taṁ

\begin{english}
con devozione rendo omaggio al Sangha.
\end{english}

Iccevam-ekanta꜕bhi꜓pūja꜕-neyya꜕kaṁ vatthuttayaṁ \\vanda꜕ya꜕tābhi꜕saṅkha꜕taṁ

\begin{english}
così è giusto offrire sincera adorazione
\end{english}

Puññaṁ ma꜕yā yaṁ ma꜕ma꜕ sabbu꜕padda꜕vā mā ho꜓ntu꜕ ve tassa꜕ pa꜕bhāva꜕-siddhi꜕yā

\begin{english}
possa il bene generato da questo buona azione dissipare ogni difficoltà.
\end{english}

Idha tathā꜓ga꜕to loke u꜕ppanno a꜕rahaṁ sammāsambuddho

\begin{english}
[Il Tathagata] è sorto in questo mondo, è il Degno, Perfettamente Risvegliato
\end{english}

Dhammo ca꜕ desi꜕to niyyāni꜕ko u꜕pa꜕sa꜕miko pa꜕rinibbāni꜕ko sa꜓mbodha꜕gāmī su꜕ga꜕tappa꜕vedi꜕to

\begin{english}
il Dhamma, che conduce oltre l'illusione, portatore di pace, la nostra guida verso il Risveglio,
Questo sentiero il Buddha ci ha indicato.

\end{english}

Ma꜓yan-taṁ dhammaṁ su꜕tvā evaṁ jānāma

\begin{english}
avendo udito l'Insegnamento, noi sappiamo che:
\end{english}

%\suttaref{DN 22}%
Jātipi꜕ dukkhā

\begin{english}
nascere è dukkha,
\end{english}

Jarāpi꜕ dukkhā

\begin{english}
 invecchiare è dukkha,
\end{english}

Ma꜕raṇampi꜕ dukkhaṁ

\begin{english}
morire è dukkha
\end{english}

So꜓ka-pa꜕rideva-dukkha꜕-domanass'u꜕pāyāsā꜓pi꜕ dukkhā

\begin{english}
tristezza, lamento e dolore, 
afflizione e disperazione sono dukkha.
\end{english}

Appiyehi꜕ sa꜓mpa꜕yogo dukkho

\begin{english}
stare con ciò che non piace è dukkha
\end{english}

Piyehi꜕ vi꜓ppa꜕yogo dukkho

\begin{english}
separarsi da ciò che piace è dukkha
\end{english}

Yamp'iccha꜓ṁ na꜕ labhati tampi꜕ dukkhaṁ

\begin{english}
non ottenere il desiderato è dukkha.
\end{english}

Sa꜓ṅkhittena pañcu꜕pādānakkha꜓ndhā dukkhā

\begin{english}
In breve i cinque aggregati soggetti d'attaccamento sono dukkha.
\end{english}

Seyya꜕thīdaṁ

\begin{english}
Essi sono:
\end{english}

Rūpūpādāna꜕kkha꜓ndho

\begin{english}
Forma,
\end{english}

Vedanūpādāna꜕kkha꜓ndho

\begin{english}
sensazione,
\end{english}

Sa꜓ññūpādāna꜕kkha꜓ndho

\begin{english}
  percezione,
\end{english}

Sa꜓ṅkhā꜓rūpādāna꜕kkha꜓ndho

\begin{english}
  formazioni mentali,
\end{english}

Viññāṇūpādāna꜕kkha꜓ndho

\begin{english}
  coscienza sensoriale.
\end{english}

%\suttaref{Trad.}%
Yesaṁ pa꜕riññāya

\begin{english}
Per la loro completa comprensione,
\end{english}

Dha꜕ramāno so꜓ bha꜕gavā evaṁ ba꜕hulaṁ sā꜓va꜕ke vi꜕neti

Evaṁ bhāgā ca꜕ panassa bha꜕gava꜕to sā꜓va꜕kesu a꜕nusā꜓sa꜕nī ba꜕hulā pa꜕vatta꜕ti

\begin{english}
il Beato, durante la sua vita insegnò spesso che:
\end{english}

%\suttaref{SN 22.90}%

Rūpaṁ a꜕niccaṁ

\begin{english}
La forma è impermanente,
\end{english}

Vedanā a꜕niccā

\begin{english}
la sensazione è impermanente,
\end{english}

Sa꜓ññā a꜕niccā

\begin{english}
la percezione è impermanente,
\end{english}

Sa꜓ṅkhā꜓rā a꜕niccā

\begin{english}
le formazioni mentali sono impermanenti,
\end{english}

Viññāṇaṁ a꜕niccaṁ

\begin{english}
la coscienza sensoriale è impermanente. 
\end{english}

Rūpaṁ a꜕nattā

\begin{english}
La forma è non-Sé,
\end{english}

Vedanā a꜕nattā

\begin{english}
la sensazione è non-Sé,
\end{english}

Sa꜓ññā a꜕nattā

\begin{english}
la percezione è non-Sé,
\end{english}

Sa꜓ṅkhā꜓rā a꜕nattā

\begin{english}
le formazioni mentali son non-Sé, 
\end{english}

Viññāṇaṁ a꜕nattā

\begin{english}
la coscienza sensoriale è non-Sé.
\end{english}

Sa꜕bbe sa꜓ṅkhā꜓rā a꜕niccā

\begin{english}
Ogni formazione è impermanente,
\end{english}

Sa꜕bbe dhammā a꜕nattā'ti

\begin{english}
Non c'è un Sé nel creato e nell'Increato.
\end{english}

%\suttaref{MN 29}%
Te ma꜓yaṁ otiṇṇāmha jāti꜕yā ja꜕rā-maraṇena

\begin{english}
Tutti noi siamo oppressi dalla nascita, dall'invecchiamento e dalla morte,
\end{english}

So꜓kehi꜕ pa꜕ridevehi꜕ dukkhe꜓hi꜕ domanassehi꜕ u꜕pāyāsehi

\begin{english}
Tristezza, lamento e dolore, afflizione e disperazione,
\end{english}

Dukkho꜓tiṇṇā dukkha꜕-pa꜕retā

\begin{english}
Oppressi da dukkha e dominati da dukkha.
\end{english}

Appeva nāmi꜓massa꜕ kevalassa꜕ dukkha-kkha꜓ndhassa꜕ anta꜕kiri꜓yā \\paññāyethā'ti

\begin{english}
Rendiamo chiara la cessazione di questa intera massa di sofferenza.
\end{english}

\begin{instruction}
  La parte seguente è recitata solo dai monaci e samanera
\end{instruction}

%\suttaref{Trad.}%
Ci꜓ra꜓-pari꜕nibbutampi꜓ taṁ bha꜕gava꜓ntaṁ uddissa a꜕raha꜓ntaṁ sammāsambuddhaṁ

\begin{english}
Ispirati dal Beato, dal Nobile, Perfettamente Risvegliato che molto tempo fa ha raggiunto il Parinibbana
\end{english}

Saddhā a꜕gārasmā anagāri꜓yaṁ pabba꜕ji꜕tā

\begin{english}
Con fede abbiamo abbracciato la vita del senza-dimora
\end{english}

Tasmi꜓ṁ bha꜕gavati brahma-ca꜕ri꜓yaṁ ca꜕rāma

\begin{english}
E come il Beato conduciamo la Vita Santa
\end{english}

Bhikkhū꜓naṁ/Sīladharā꜓naṁ si꜓kkhāsā꜕jīva꜕-samāpannā

\begin{english}
Avendo intrapreso pienamente l'addestramento monastico
\end{english}

Taṁ no brahma-ca꜕ri꜓yaṁ imassa꜕ kevalassa꜕ dukkha-kkha꜓ndhassa꜕ anta꜕kiri꜓yāya sa꜓ṁva꜓tta꜕tu

\begin{english}
Possa questa Vita condurci alla fine di questa intera massa di sofferenza.
\end{english}

\clearpage

\begin{instruction}
  La parte seguente è recitata solo dai laici
\end{instruction}

Ci꜓ra꜓-pari꜕nibbutampi꜓ taṁ bha꜕gava꜓ntaṁ saraṇaṁ ga꜕tā

\begin{english}
Il Beato che molto tempo fa ha raggiunto il Parinibbana è il nostro rifugio
\end{english}

Dha꜓mmañca sa꜓ṅghañca

\begin{english}
Così anche il Dhamma e il Sangha.
\end{english}

Tassa bha꜕gavato sā꜓sanaṁ yathā꜓-sati yathā꜓-balaṁ manasika꜕roma a꜕nupaṭipa꜓jjāma

\begin{english}
Con presenza mentale e determinazione pratichiamo il suo Insegnamento
\end{english}

Sā꜓ sā꜓ no pa꜕ṭi꜓patti

\begin{english}
Possa quindi questa via
\end{english}

Imassa꜕ kevalassa꜕ dukkha-kkha꜓ndhassa꜕ anta꜕kiri꜓yāya sa꜓ṁva꜓tta꜕tu

\begin{english}
Condurci alla fine di questa intera massa di sofferenza.
\end{english}

\clearpage

\chapter{Omaggio finale}

%\suttaref{Trad.}%
[Arahaṁ] sammāsambuddho bha꜕gavā

\begin{english}
Il Beato, il Nobile Perfettamente Risvegliato
\end{english}

Buddhaṁ bha꜕gavantaṁ a꜕bhi꜓vādemi

\begin{english}
Rendo omaggio al Buddha, al Beato.
  \instr{Inchino}
\end{english}

[Svākkhā꜓to] bha꜕gava꜓tā dhammo

\begin{english}
L'insegnamento da lui completamente esposto,
\end{english}

Dhammaṁ namassāmi

\begin{english}
Mi inchino al Dhamma.
  \instr{Inchino}
\end{english}


[Supaṭi꜕panno] bha꜕gava꜕to sāvaka꜕saṅgho

\begin{english}
I discepoli del Beato, che hanno ben praticato,
\end{english}

Sa꜓ṅghaṁ na꜕māmi

\begin{english}
Mi inchino al Sangha.
  \instr{Inchino}
\end{english}

% End of morning-chanting.tex
